\chapter{Anatomy of a Modern Fuzzer}
\label{ch:afl}

While the components outlined in Section~\ref{sec:ff} establish the blueprint for most feedback-driven fuzzers, translating this architecture into a performant tool capable of operating at an industrial scale requires highly specific engineering implementations. The cycle of scheduling inputs, executing the target and evaluating feedback demands heavy performance optimisations to maintain efficiency in real-world usage.

AFL++ \cite{fioraldi2020afl++} has established itself as the \textit{de facto} standard in modern coverage-guided grey-box fuzzing. By consolidating diverse research advancements (such as \cite{aschermann2020ijon}, \cite{aschermann2019redqueen} and \cite{bohme2016coverage}) into a single project, it serves as the premier baseline for contemporary vulnerability research. Understanding the architectural choices provides context that is fundamental to identifying constraints that limit traditional edge coverage metrics.

\section{Compilation and Instrumentation}
\label{sec:compandinst}

Without execution telemetry, a fuzzer is essentially blind. To achieve visibility, the fuzzer must inject monitoring logic into the target program. Implementing an observer generally involves modifying the binary to report every basic block transition traversed during the execution via instrumentation.

However, instrumentation is not a monolithic process. In AFL++, there are multiple compile-time strategies available, which heavily dictate both the accuracy of the resulting coverage map and the performance overhead of the fuzzing loop.

AFL++ introduces three distinct compilation modes:
\begin{itemize}
    \item \textit{LLVM Mode}: Operates at the LLVM IR level, removing the need for assembly rewriting. It provides efficient basic-block instrumentation while laying the groundwork for more advanced features.
    
    \item \textit{LTO Mode}: Uses Link-Time Optimization to guarantee a collision-free coverage bitmap.

    \item \textit{GCC Plugin Mode}: It uses GCC compiler plugins, since not all codebases can be easily migrated to the LLVM infrastructure. It provides an essential, high-performance bridge without requiring a compiler toolchain overhaul.
\end{itemize}

In scenarios where source code is unavailable (such as proprietary or legacy software) AFL++ leverages emulation or dynamic binary instrumentation --- such as QEMU or FRIDA --- to inject execution tracking logic at runtime, accepting higher performance overhead for maintaining visibility.

\subsection{LLVM Mode}
\label{sub:llvmmode}

LLVM mode --- typically invoked using the flag \texttt{--afl-llvm}, the environment variable \texttt{AFL\_CC\_COMPILER=LLVM}, or via the compiler wrappers \texttt{afl-clang-fast} and \texttt{afl-clang-fast++} --- serves as the modern standard for instrumenting source-available targets.

This mode operates strictly at the LLVM Intermediate Representation (IR) level, removing the need for brittle assembly rewriting. Consequently, this allows the compiler to treat injected instrumentation as native code, allowing standard optimisation passes to optimise it alongside the target program and reduce performance overhead. Furthermore, operating at the IR level also means that the instrumentation is architecture-agnostic and generally independent of the target CPU architecture.

To fully leverage the LLVM framework, AFL++ implements its instrumentation as a custom LLVM compiler pass. Passes execute all transformations and optimisations of the compiler, compute the analysis results required by these operations, and serve as the architectural framework for compiler code \cite{writingllvmpass}.

The compiler wrappers work as drop-in replacements for the standard C/C++ compilers. When invoked, they execute the standard Clang compiler but instruct it to load AFL++'s custom shared object --- the instrumentation pass --- during the intermediate compilation pipeline. The final linking phase is equally important, as the compiler wrapper links the static library for the compiler runtime, whose primary responsibility is setting up the shared memory region for the coverage map and initializing the forkserver. 

\subsubsection{PCGUARD Instrumentation}

AFL++ heavily leverages and extends LLVM's \texttt{SanitizerCoverage} infrastructure. Rather than manually traversing the CFG to find every edge, AFL++'s \texttt{SanitizerCoveragePCGUARD.so.cc} pass performs the following:
\begin{enumerate}
    \item \textit{Block selection and pruning:} The pass iterates over all basic block structures in a function, then uses Dominator and Post-Dominator trees to identify blocks where coverage tracking is redundant (e.g., a block that has a single predecessor and a single successor). These blocks can be skipped \cite{agrawal1994dominators} to reduce overhead.

    \item \textit{Guard Allocation:} For every instrumented block, the compiler reserves a 32 bit slot in a dedicated ELF section (\texttt{\_\_sancov\_guards}). At compile-time, these are uninitialized integers. 

    \item \textit{Inline IR Injection:} For each selected basic block, LLVM IR is injected at the start of the block:
    \begin{itemize}
        \item It creates an \texttt{IntToPtr} lookup into the \texttt{sancov\_guards} array using a unique index for that block.
        
        \item It loads the 32-bit coverage ID from that guard slot.
        
        \item It uses a \texttt{GetElementPtr} (GEP) instruction to calculate the exact byte address within the \texttt{\_\_afl\_area\_ptr} (the shared memory map bridging the fuzzer and target) using the loaded coverage ID.
        
        \item It loads the current 8-bit counter from that shared memory byte, increments it, and stores it back.
    \end{itemize}
    See Appendix \ref{app:llvmir} for an example of standard injected LLVM IR.
\end{enumerate}

\paragraph{Initialisation of Guards} The mode relies on an array of guard variables, which must be initialised at runtime. The compiler groups all \texttt{\_\_sancov\_guards} across different translation units. At load time, a constructor function provided by the compiler runtime is executed (\texttt{\_\_sanitizer\_cov\_trace\_pc\_guard\_init}). This runtime function iterates over all 32-bit guard pointers, assigning a random pseudo-unique ID to each one, bounded by the shared memory size. This allows mapping the CFG basic blocks to random slots in the fuzzer's shared bitmap.

\subsubsection{Classic Tracing}

When PCGUARD is not available, AFL++ falls back to "Classic" tracing mode. The pass \texttt{afl-llvm-pass.so.cc} computes edge coverage using the original AFL coverage formula \cite{zalewski2014afltech}:
\[
    \text{map}[\text{cur\_loc} \oplus (\text{prev\_loc} \gg 1)] \mathrel{{+}{=}} 1
\]

In this design, \texttt{cur\_loc} is a pseudo-random identifier statically assigned at compile-time, while \texttt{prev\_loc} holds the identifier of the previously executed block. To maintain execution context across the program, the injected instrumentation must save the current block's identifier into a thread-local variable before transferring control flow to the subsequent block.

The bitwise right shift ($\gg 1$) is necessary to distinguish edge $A \rightarrow B$ from $B \rightarrow A$, since the XOR operation is commutative. Not modifying one of the operands would result in symmetric transitions, losing the sequence's directionality. It also prevents self-loops ($A \rightarrow A$) from always being mapped to index 0.

\subsection{LTO Mode}

In standard LLVM mode, every C/C++ source file (translation unit) is compiled and instrumented entirely independently, with no view of the global program. This inevitably leads to map collisions when linking multiple object files, as random basic block IDs might overlap.

Link-Time Optimization (LTO) mode addresses this inefficiency by shifting the instrumentation process to the final linking phase. By operating with global visibility, the AFL++ LTO instrumentation pass (\texttt{SanitizerCoverageLTO.so.cc}) can uniquely identify each edge. As a result, LTO mode guarantees a collision-free coverage bitmap. By eliminating index overlaps, the fuzzer receives better edge coverage granularity, improving the accuracy of the feedback loop.

However, these benefits are tied to stricter toolchain requirements, requiring LLVM version 12+ to natively allow the LLVM linker to load plugins.

\paragraph{The LTO Solution} With this mode, AFL++ defers final compilation and optimization until the linking phase, intercepted via the compiler wrappers \texttt{afl-clang-lto} and \texttt{afl-clang-lto++}.

When compiling source code with the wrappers, the compiler passes the \texttt{-flto=full} flag to LLVM, instructing it to emit LLVM bitcode instead of the final machine code.

At link time, the custom linker \texttt{afl-ld-lto} invokes AFL++'s specialized linker pass. With global visibility into the CFG of the entire program, there is no need for random IDs, instead, it iterates through every block across the entire binary and assigns a contiguous, monotonically increasing identifier to every edge.

\subsubsection{LTO-Exclusive Features}

\paragraph{Autodictionary} By analysing the whole program at once, the linker can see every single constant string across all linked libraries and object files. During the LTO pass, it automatically scrapes these strings and generates a binary dictionary structure explicitly embedded within a specific ELF section of the target binary.

The AFL++ runtime extracts this dictionary when starting the target, immediately integrating it into its mutation strategies.

\paragraph{Fixed Memory Map Address} To eliminate the overhead associated with dynamically looking up the shared memory map pointer in the Global Offset Table (GOT), LTO mode can hardcode the absolute memory address of the shared memory region directly into the assembly of the coverage instructions. 

This can crash targets that utilize early constructors, deferred forkservers or anything that might conflict with OS memory layout or ASLR.

\paragraph{Document Edge IDs} The compiler outputs a comprehensive mapping of every assigned Edge ID to its corresponding function and file name.

\subsection{GCC Plugin Mode}

LLVM and LTO modes inherently rely on the target software being compatible with the LLVM compiler toolchain. In practice, not all codebases can be easily migrated to the LLVM infrastructure. Legacy systems and heavily architecture-dependent projects often cannot be compiled using Clang.

To maintain efficient execution insights, AFL++ incorporates the GCC plugin mode, replacing the assembly-level rewriting approach of the original \texttt{afl-gcc}. Instead of inserting assembly blocks directly into the compiled output, this mode operates at the compiler level via true GCC plugins. This ensures that the target can still benefit from compile-time instrumentation and optimisations while remaining tied to the GCC ecosystem.

To build with \texttt{afl-gcc-fast}, the host system must have GCC version 4.5.0 or newer and the corresponding \texttt{gcc-VERSION-plugin-dev} installed.

\paragraph{GIMPLE Pass Injection} The GCC compiler uses several intermediate representations, most notably GIMPLE \cite{gccinternals}. The AFL++ plugin registers a custom pass that operates on the GIMPLE trees. The pass \texttt{afl-gcc-pass.so.cc} acts as a GIMPLE pass which iterates over all basic blocks within a function and: 
\begin{enumerate}
    \item Assigns a random integer ID to all basic blocks.

    \item Constructs GIMPLE statements to calculate the standard AFL edge coverage formula.

    \item Emits a GIMPLE memory load/store sequence to increment the corresponding byte in the \texttt{\_\_afl\_area\_ptr} shared memory map.
\end{enumerate}

Because it relies strictly on GCC's plugin API and GIMPLE, this mode lacks the deeper analysis features found in LLVM, such as LAF-Intel (Split Compares), autodictionary and context-sensitive coverage.

\subsection{Black-Box Binary Instrumentation}
\label{subsec:bbb-instr}

When source-code is unavailable, compile-time instrumentation is impossible. To fuzz a black-box binary efficiently, dynamic binary instrumentation (DBI) or static binary rewriting is required. 

AFL++ offers extensive support for binary-only targets through several modes and integrations. The choice of mode depends largely on architecture, operating system, performance requirements and stability of the target. 

FRIDA or QEMU mode with persistent mode are the fastest, provided stability is high enough.

\subsubsection{QEMU Mode}

QEMU mode utilizes a patched version of QEMU running in user space emulation to instrument black-box, closed source binaries dynamically. It translates the target binary instructions into Tiny Code Generator (TCG) intermediate representation and injects AFL++ coverage collection instrumentation.

AFL++ can hook directly in the translation phase from machine code into TCG. When a new basic block is encountered and translated, the equivalent of the AFL coverage logging instruction is injected into the TCG block. Once translated, these cached blocks execute natively on the host.

\paragraph{Persistent Mode} It's the most significant performance enhancement, speeding up execution by several factors. This mode allows a target loop to be executed continuously within the same process, without forking for every attempt.

\subsubsection{FRIDA Mode}

This mode is an alternative to QEMU mode, usually with comparable performance, if not slightly faster. It is designed to be a drop-in replacement for QEMU mode, mimicking its environment variables (replacing \texttt{QEMU} with \texttt{FRIDA} in the name).

FRIDA mode operates by injecting a shared library \texttt{afl-frida-trace.so} into the target process. This relies on the \texttt{frida-gum} devkit and its dynamic translation engine ``\textit{Stalker}''.

Instead of full emulation, Stalker dynamically translates code basic block by basic block, injecting assembly snippets directly into the translated blocks. This model allows the fuzzer to inject execution tracking logic directly into the memory space of a running process.

\paragraph{FASAN: FRIDA Address Sanitizer Mode} Unlike QEMU mode --- which crashes if combined with ASAN --- FASAN allows injecting the same memory safety constraints at runtime. The dynamic loader intercepts standard memory functions and FRIDA dynamically instruments all memory operand instructions in the target binary, inserting calls provided by the ASAN dynamic shared object to validate read/writes operations against the shadow memory.

\subsubsection{Other Modes}

\paragraph{Nyx Mode} This mode shifts the paradigm to full-system emulation. Built upon a modified KVM and QEMU hypervisor, Nyx is designed to handle complex, stateful targets, kernel components or targets that cannot utilize persistent mode. Instead of utilizing forks, it utilizes KVM-based snapshots to revert the VM state.

\paragraph{Unicorn Mode} It's based on the Unicorn Engine, a fork of QEMU, and the instrumentation is, therefore, very similar. While QEMU provides a full Linux user-space environment, Unicorn Mode provides no environment at all; it simply emulates CPU instructions. Consequently, this mode requires writing a custom runtime harness to handle all memory mapping, register initialization, and syscall interception manually.

\paragraph{CoreSight Mode} This is AFL++'s hardware-assisted tracing solution, specifically designed for Arm64 architectures. It leverages ARM CoreSight technology to use the physical CPU's hardware trace macrocells to record control flow execution \cite{shan2023crowbar}. The CPU directly writes execution branch data into memory, which then has to be decoded. 

\paragraph{Binary Rewriters} These are external tools that can inject AFL++ coverage instrumentation directly into the binary either statically on disk --- such as ZAFL \cite{nagy2021breaking}, RetroWrite \cite{dinesh2020retrowrite} or Mcsema \cite{dinaburg2014mcsema} --- or dynamically at load/run time --- such as Dyninst \cite{buck2000api}.

\section{Runtime Components}

Compilation and instrumentation equip the target binary with the tools to report execution paths, but once visibility into the target is established, the challenge of CGF shifts to execution throughput. Processing millions of mutated inputs via standard system execution flows can introduce heavy performance bottlenecks.

The fuzzer requires an execution environment capable of managing the fuzzing loop, delivering mutated inputs, and evaluating the target's behaviour at scale. To achieve efficiency, modern fuzzers use high-speed inter-process communication (IPC) and specialized state management.

\subsection{The Forkserver Architecture}

The normal system execution flow for executing a program is to call \texttt{fork} and \texttt{execve}, but the latter is extremely slow, as it requires the kernel to load the binary, set up the memory layout, and invoke the dynamic linker. 

The \textit{forkserver} avoids this overhead. The fuzzer executes the target program only once. The target program initializes itself and stops before the main fuzzing loop, either automatically at \texttt{main} or via \texttt{\_\_AFL\_INIT()}. At this point, it acts as a "forkserver": for every fuzzing iteration, the fuzzer sends a signal to the forkserver, which simply calls \texttt{fork}, letting the newly created child process inherit the already-initialized memory state. The child then executes the target function with the new input and exits. The forkserver parent waits for the child and reports its exit status back to the fuzzer.

\paragraph{Communication Channels} Fuzzer and forkserver communicate via two standard pipes: 
\begin{itemize}
    \item A \textit{Control} pipe \texttt{FORKSRV\_FD}, used by the fuzzer to send commands to the forkserver.

    \item A \textit{Status} pipe \texttt{FORKSRV\_FD + 1}, used by the forkserver to send status updates back to the fuzzer.
\end{itemize}

\paragraph{Handshake protocol} AFL++ introduces a protocol to allow fuzzer and target to negotiate capabilities dynamically at startup.

It includes an initial handshake to check the supported forkserver version; then, the fuzzer and target negotiate advanced features through a 32-bit bitmask.

\paragraph{Execution Loop} Once the handshake is complete, the standard execution loop begins for each input generated by the fuzzer: 
\begin{enumerate}
    \item The fuzzer writes a 4-byte command to the control pipe, typically instructing it to start.

    \item The forkserver reads the command. If the pipe is closed, the fuzzer exited and the forkserver terminates itself. Otherwise, it calls \texttt{fork}, lets the child run the target logic with the new input, and writes the PID of the child process to the status pipe.

    \item The forkserver then waits for the child to finish and writes the exit code to the status pipe.

    \item The fuzzer reads the status, updates the coverage bitmap, calculates the feedback and repeats the cycle.
\end{enumerate}

\paragraph{Alternative Backends} The forkserver architecture allows various execution backends, such as the \textit{fauxserver} --- which uses \texttt{execv} for targets incompatible with the fork architecture --- or other modes such as the ones described in Section \ref{subsec:bbb-instr}.

\subsubsection{Persistent Mode}

Even with a forkserver, calling \texttt{fork} thousands of times per second incurs significant overhead. Persistent mode solves this by doing a single \texttt{fork} and inserting a loop around the target function. The child then repeatedly reads the fuzzer input, executes the target function and resets any necessary state, without ever exiting. 

When the loop counter reaches its limit, the child exits and a new one is created from the forkserver. This limits the impact of memory leaks while still achieving superior performance.

AFL++ includes features for tracking persistent mode executions. Because the target state might not be perfectly reset between iterations in persistent mode, a crash might depend on the specific sequence of inputs executed before the one that causes the crash. With \texttt{AFL\_PERSISTENT\_RECORD} enabled, the forkserver tracks the last $n$ inputs processed by the child, allowing the fuzzer to replay the exact state transitions that led to the crash.

\subsection{Shared Memory}
\label{sub:shm}

The primary mechanism used to orchestrate high-throughput communication is shared memory. By mapping common memory regions into the address spaces of both fuzzer and target, AFL++ avoids the latency overhead associated with traditional IPC --- such as using pipes or file system I/O.

The shared memory architecture in AFL++ relies on zero-copy data exchange. When the fuzzer is initiated, it allocates memory regions for feedback data and operational purposes. Upon execution, the target binary retrieves, via environment variables, the identifiers for these pre-allocated regions, attaching them to its own virtual address space.

\subsubsection{Lifecycle and Allocation Mechanisms}

The management of the shared memory is mainly governed by a single module (\texttt{afl-sharedmem.c}) which abstracts the OS paradigms into an API, defined in the corresponding header file. There is a data structure maintaining everything related to the state, file descriptors, object identifiers and virtual memory pointers for the shared regions.

The initialisation routine handles allocating and preparing the necessary memory segments, as well as constructing the corresponding environment strings. Along with the primary coverage map, it handles the creation of auxiliary maps, when enabled (e.g., CmpLog, ASan). A deinit function provides the cleanup counterpart, unmapping virtual memory and unsetting all environment variables.

\paragraph{Large Page Allocation} To minimize Translation Lookaside Buffer (TLB) misses during execution loops, AFL++ attempts to utilize \textit{large pages}. This decreases memory management overhead and increases performance.

\subsubsection{Coverage Map}

The primary coverage map acts as the central data structure for evaluating execution feedback. Instantiated within the shared memory region during the initialisation phase, it is conventionally structured as a fixed-size (although configurable) contiguous byte array.

After the forkserver signals the termination of a child process, the fuzzer directly analyses the modified map without further IPC. Each byte in the array corresponds to the execution frequency of a specific control-flow edge. When a transition from basic block $A$ to $B$ occurs, the instrumentation logs this edge by incrementing a byte in the map.

The coverage map's lifecycle is split between two processes: the fuzzer --- which creates and analyses it --- and the target --- which attaches to and mutates it.

\paragraph{Post-Execution Analysis} After the target child process terminates, control returns to the parent process, which must efficiently analyse the resulting map to determine if the execution yielded interesting behaviour. 

The fuzzer first checks for novelty in the generated map by comparing it against a persistent \texttt{virgin\_bits} map, which tracks all previously-discovered edges and is updated when novel transitions are discovered.  

To combat path explosion caused by trivial variations in the number of loop iterations (e.g., observing an edge 7 times instead of 6), the fuzzer groups execution counts into one of 8 pre-defined buckets. Changes within the same bucket are ignored, mitigating state explosion.

The input is considered "interesting" if it returns new coverage or a new count for one of the buckets. The input is subsequently saved to the queue directory and becomes a seed for future mutation cycles.

\subsection{Queue and Scheduling}
\label{subsec:qandsched}

In AFL++, the queue is not a simple data structure. It serves as the evolving repository for the fuzzer's seed corpus. Every input that triggers new coverage or unique behaviour is preserved as a distinct queue entry, serving as material for future mutations. It is the queue's responsibility to orchestrate the fuzzing loop by deciding which entry to schedule next and how much \textit{energy} (i.e., fuzzing time) to allocate to it.

Before being actually inducted into the queue, the input undergoes a calibration phase. The test is executed multiple times to ensure the observed behaviour is deterministic rather than an anomaly.

Furthermore, there is an array that keeps the best queue entry --- in terms of speed and size --- for each byte of the coverage map. Only the best entry for each edge is kept.

\paragraph{Scoring} When selecting a test case from the queue to undergo mutation, AFL++ must determine the intensity and duration of the fuzzing cycle. This is codified as the \texttt{perf\_score}. The calculation of the score is a composition of several heuristics, all modulated by a \textit{power schedule}.

The base metrics taken into account are: 
\begin{itemize}
    \item \textit{Execution speed:} Faster inputs are prioritized, as they're easier to fuzz.

    \item \textit{Bitmap density:} Entries that traverse more unique edges --- i.e., have a fuller bitmap --- are assumed to be more valuable.

    \item \textit{Temporal handicap:} When a new queue entry is discovered late in the fuzzing campaign, it has inherently missed the previous mutation cycles that other entries have enjoyed; to mitigate this, newly discovered entries have a higher multiplier in scoring.

    \item \textit{Generational Depth:} The distance, in number of generations, from the starting corpus is treated as a proxy for semantic complexity; inputs with higher depth are given higher multipliers.
\end{itemize}

There are several power schedules, used as "macro-strategies" that fundamentally alter the calculated score, in one way or another.

\paragraph{Selection} AFL++ does not simply iterate over the queue sequentially. It uses weighted random selection via an alias table.

Each entry's weight is derived from its \texttt{perf\_score} and the next queue entry is selected by using the alias table. This ensures that \texttt{favored} and highly scored entries are selected far more frequently than redundant or slow entries.

\paragraph{Culling} To prevent the fuzzer from wasting time on redundant paths, a greedy algorithm is called periodically to calculate a minimal set of test cases that hit every known edge. These are marked as \texttt{favored}. During the main fuzzing loop, there is an extreme bias towards the \texttt{favored} entries.

Efficient entries often cover large portions of code; selecting a single entry from all the top-rated ones might clear thousands of bits from the queue. This ensures that the final set of favoured entries remains much smaller than the total queue.

\subsection{Compiler Runtime}

The instrumentation logic cannot work without the AFL++ compiler runtime (\texttt{afl-compiler-rt.o.c}). The runtime is responsible for bridging the gap between the instrumented edges embedded in the target binary and the fuzzer process. It is linked as a static object, and as such constructors inside the runtime are executed before the target main function is invoked.

The instrumentation process produces code that references symbols provided by the runtime, including generic variables such as the ones used for the shared memory.

The runtime operates primarily as an initialisation bootstrap and execution harness. It intercepts the normal execution flow of the target program to establish the shared memory region, initialise the forkserver, manage persistent mode, and support advanced features.

\paragraph{Shared Memory} AFL++ introduced dynamic map sizing to support larger targets without excessive map collisions. 

The instrumentation passes export a symbol representing the total number of instrumented edges. The runtime evaluates the value and, if larger than the default size, the runtime sets the size of the map to accommodate the exact number of edges.

Depending on the platform and configuration, the runtime attaches the shared memory using either POSIX shared memory or System V shared memory. The base address of the mapped memory is assigned to the global variable \texttt{\_\_afl\_area\_ptr}.

\paragraph{Forkserver} Inside the runtime resides the forkserver initialization hook, used to pause the process and establish the communication channel with the fuzzer process. 

Following the communication handshake, instead of terminating after a single input, the forkserver enters an infinite loop of:
\begin{itemize}
    \item Waiting for a command from the fuzzer.

    \item Cloning itself, executing program logic in the child.

    \item Waiting for the end of execution and reporting the status to the fuzzer process.
\end{itemize}

\paragraph{Deferred Initialisation} By default, AFL++ stops the forkserver just before the main function of the target process. While safe, this means that every cloned child must re-execute the entirety of the initialisation code. If a target performs expensive one-time setups --- e.g., parsing large configuration files --- fuzzing performance will be severely degraded.

Deferred initialisation solves this by allowing the forkserver to start at a point chosen by the user, where the expensive setup is completed.

In this case, the compiler runtime skips calling the forkserver constructor and instead calls the manual forkserver initialisation at the specified time. 

\subsection{Mutation Engine}

AFL++'s mutation engine is the component primarily responsible for input generation. Unlike blind random fuzzing, AFL++ employs a more sophisticated, coverage-guided evolutionary algorithm. It utilizes a structured and heuristic-driven pipeline to alter seed inputs. By continually adapting its strategies based on execution tracing, the engine ensures the computational effort is focused on generating the best possible inputs.

\paragraph{Phases} The standard execution of the default mutation engine for any given input follows a sequential progression:
\begin{enumerate}
    \item Calibration and Trimming: Before mutation, the input may be calibrated to confirm the execution trace and trimmed to reduce size without losing coverage.

    \item Deterministic stages: if enabled and the input has not already been subjected to them, AFL++ applies the following stages: bitflips, general arithmetic, insertion of common interesting values and dictionary insertions.

    \item Non-Deterministic Stage (Havoc): multiple mutations stacked on top of each other in a highly randomized fashion. The number of cycles depends on the seed's performance score and queue depth; favourable inputs receive significantly more havoc time than uninteresting ones.

    \item Splicing: if the havoc stage fails, the engine will attempt splicing an input with another one from the queue.
\end{enumerate}

The early stages guarantee that every single byte will be modified sequentially. This is guaranteed to find simple syntactic branches but can be computationally expensive. For these reasons, it is never repeated for the same queue entry once completed.

\section{Advanced Features}

Standard fuzzing campaigns often encounter significant bottlenecks when faced with highly structured input formats, deep execution states, or complex checksums. When applied to these targets, fuzzers can run into "path starvation", expending large amounts of energy on superficial code paths.

To overcome problems such as these, AFL++ integrates mechanisms to bypass complex program states and modify the behaviour of mutation and scheduling.

\subsection{Advanced Coverage Metrics}

A fundamental constraint in CGF efficiency is how the fuzzer perceives the program state. Standard edge-based feedback allows the fuzzer to observe only the immediate transition between two basic blocks, sometimes resulting in lost feedback if novel program states reuse code paths.

AFL++ implements different coverage metrics that attempt to solve this problem by enriching the coverage map calculation with additional information, allowing the fuzzer to better differentiate between execution paths.

\subsubsection{Context-Sensitive Branch Coverage}

In some codebases, a single utility function might be called from many different locations. Since the standard coverage metric is based on a single edge, to the fuzzer the internal control flow of the utility function appears to be identical, regardless of the component that called it. To solve this insensitivity to context, the fuzzer incorporates the state of the call stack into the edge index computation. 

AFL++ does this by assigning a unique ID to every function at compile time and saving a hash of the call stack inside Thread-Local Storage (TLS), which is then XORed with the standard edge index calculation. 

The hash is calculated by injecting an instruction to XOR the TLS value with the function ID at each function start and return.

\subsubsection{N-Gram Branch Coverage}

Traditional grey-box fuzzing measures code coverage considering only the current basic block and the preceding one. While efficient, this can lead to \textit{path collision}, where semantically divergent execution paths can converge on the same block transition, indistinguishable to the fuzzer.

To solve this, AFL++ integrates N-Gram Branch Coverage \cite{wang2019sensitive}. This metric expands the state memory of the fuzzer, allowing it to trace execution histories of length $N$.

\begin{figure}
    \centering
    \resizebox{0.5\textwidth}{!}{
        \begin{tikzpicture}[
    scale=0.5,
    ->, % makes edges directed
    >=stealth, % arrow head style
    node distance=0.9cm and 1.5cm, % vertical and horizontal spacing
    thick,
    % Styling for the nodes
    cfgNode/.style={
        rectangle, 
        rounded corners, 
        draw, 
        minimum width=1.618cm, 
        minimum height=1cm, 
        font=\sffamily\bfseries
    }
]

    % Define Nodes
    \node[cfgNode] (A) {$A$};
    \node[cfgNode] (X) [above left=of A] {$X$};
    \node[cfgNode] (Y) [below left=of A] {$Y$};
    \node[cfgNode] (B) [right=of A] {$B$};

    % Draw Edges
    \draw (X) -- (A);
    \draw (Y) -- (A);
    \draw (A) -- (B);

\end{tikzpicture}
    }
    \caption{Example of CFG}
    \label{fig:cfgex1}
\end{figure}

In the example represented in Figure \ref{fig:cfgex1}, under standard AFL++, the transition $A \rightarrow B$ looks identical regardless of the starting block. Under an N-Gram model with $N > 2$, the two possible paths evaluate to distinct hashes, rewarding the fuzzer for discovering a new path.

AFL++ implements this by keeping track of the last $N-1$ previous blocks via an array and injecting instrumentation to modify it at every basic block.

This mode substantially enhances the fuzzer's sensitivity to complex state machines and sequential parsing logic, but it also greatly increases the number of unique identifiers generated.

\subsubsection{Ball-Larus Path Coverage}

Each edge, even when augmented via context sensitivity or N-Gram coverage, inherently remains a decomposition of control flow. However, developers often think about code in terms of logical paths. 

To align the fuzzer's feedback with the true semantic topology of the target program, Ball-Larus Path Coverage treats functions as Directed Acyclic Graphs (DAGs), assigning a unique identifier to every complete acyclic route through a function. 

The Ball-Larus algorithm \cite{ball1996efficient} establishes a method for generating a perfect minimal hash for all paths in an acyclic control graph.

The CFG of a function is transformed into a DAG by removing all loop-back edges. Loops are effectively treated as a single block. Then the graph is traversed from exit back to entry and an integer is assigned to each edge such that the sum of the weights along any path from entry to exit uniquely equates to an integer in the range $[0, NPaths - 1]$ --- with $NPaths$ being the total number of paths in the graph.

During execution, an integer is set to 0 on function entry. As control flows through edges, the weight of each traversed edge is added to the counter. On function exit, the final value of the accumulator uniquely represents the path traversed.

Because of the possible combinatorial explosion in the number of total acyclic paths, AFL++ introduces collapse heuristics. There are various configurations, from strict adherence to the Ball-Larus algorithm to considering only behavioural path divergence, tracked by using a maximum instead of a sum for the number of path counts. 

\subsection{CmpLog Mode}

Standard coverage-guided fuzzers rely on byte-level bit flips and other random mutations to explore an application's state space. While these mutations are effective for broad structural vulnerabilities, they struggle when encountering blocks such as multi-byte comparisons, checksums or parser constraints.

When employing random mutations, the probability of correctly guessing a 64-bit integer is $1$ in $2^{64}$. Solving these constraints has traditionally been done with analysis techniques such as symbolic execution. These approaches, while correct, can be orders of magnitude slower than native execution.

To bypass the costs of symbolic execution, the Redqueen algorithm \cite{aschermann2019redqueen} proposed a solution rooted in the concept of input-to-state correspondence, i.e., the premise that during the execution of a program parts of the input frequently end up in the program state unmodified, or with predictable encoding.

When a highly constrained conditional branch is evaluated, the operands of the comparison instruction or routine often consist of a component derived from the fuzzer's input, an expected constant, or a dynamically generated value.

If a fuzzer can actively monitor and log these operands during execution, it can avoid formally modelling the application's constraints. The fuzzer can directly observe that a specific value is being compared against the execution context. By actively substituting the expected value into the original input the constraint is immediately solved.

AFL++ realises this framework through its CmpLog instrumentation mode. The fuzzing process is divided into two binaries: one for coverage and one with specialized hooks before every memory or string comparison.

When the fuzzer discovers a new path that is heavily constrained or stalling, the fuzzer submits the input to the CmpLog binary. The fuzzer then parses the information given by the additional instrumentation, identifying the values required to traverse previously failed conditional branches.

The Redqueen mutator then attempts to localize these values within the original input --- a process known as \textit{Colourisation}. If the fuzzer can determine which input bytes directly influence comparison operands, it injects the logged target values to satisfy the constraints.

\subsection{MOpt-Adaptive}

Traditionally, stochastic mutations in fuzzers are determined by static probability distributions. This does not take into account the responses of different target applications. For instance, a parser might benefit disproportionately from block splicing, while a cryptographic library might yield more coverage through bit flips.

The MOpt algorithm \cite{lyu2019mopt} models the selection of mutation operators as a Multi-Armed Bandit problem (MAB). The initial implementation utilized Particle Swarm Optimization (PSO) to dynamically update the selection probabilities of mutation operators based on their empirical efficiency during a "pilot" phase.

AFL++ evolved this concept, abandoning PSO in favour of a context-aware MAB framework, stabilised by simulated annealing, exponential decay, and an exploration-exploitation floor. The innovation lies in recognizing that the efficacy of a mutation is not strictly a property of the target application, but rather of the fuzzing context. 

\paragraph{Contextual space} The algorithm partitions the fuzzer's state into a set of discrete contexts, parametrised by two dimensions: 
\begin{itemize}
    \item \textit{Input mode:} Categorising the structure of the seed input; it can be "Generic", "Text" and "Binary".

    \item \textit{Fuzzing mode:} Categorising the objective of the current fuzzing iteration. This can be "Explore", aimed at broad state discovery, and "Exploit", aimed at highly favoured paths.
\end{itemize}
The probability distribution of mutation operators is independently learned and maintained for each context.

\paragraph{Efficacy Measurement} To establish the utility of each mutation operator $i$, MOpt-Adaptive tracks two primary metrics per operator within the active context: "Uses" $U_{i,t}$, representing the number of times an operator has been applied, and "Finds" $F_{i,t}$ representing the number of times the operator successfully contributed to the discovery of a new path.

Both metrics are subject to exponential decay, to prevent historical performance from dominating the future utility. This way, the learned probability distribution remains sensitive to recent features of the application.

There is also a starting static distribution $P$, whose weight decreases linearly over time, down to a minimum floor.

The final probability weight $W_i$ for an operator is a convex combination
\[
    W_i = (1 - \varepsilon) \cdot \frac{F_i}{U_i} + \varepsilon \cdot P_i
\]

This weight is quantised in discrete slots, for faster look-up by using a table. There is an exploration floor; no operator can have fewer than one slot.

The learned distribution is also compared against the predefined static distribution. If the learned array fails to demonstrably outperform the static array's rate of discovery, the fuzzer will naturally fall back to the static distribution.

\subsection{Custom Mutators}

To breach parsing boundaries in complex targets, the mutation engine must have awareness of the input domain. By exposing the custom mutator API, AFL++ allows injecting domain-specific knowledge into the fuzzing loop.

A custom mutator is an external module --- dynamically loaded into the AFL++ process --- that intercepts the mutation phase, allowing for structure-aware alterations. The custom mutator is positioned as the first non-deterministic stage, ensuring that domain-specific mutations are prioritized. It is also possible to enforce the use of custom mutators exclusively, disabling AFL++'s intrinsic mutation stages. This capability is achieved through a series of callbacks managed by the fuzzer process.

By implementing hooks such as \texttt{afl\_custom\_fuzz} and \texttt{afl\_custom\_post\_process}, mutators can translate inputs into intermediate representations, perform targeted structural manipulations (e.g., altering Abstract Syntax Tree nodes), and serialize them back into a valid format. When confronting targets with syntactic constraints, this domain-aware approach ensures that test cases bypass superficial validation layers, allowing the evolutionary algorithm to focus on uncovering vulnerabilities within the deeper logic of the application.
